% 本文件独立定义所有所需局部写法；由主稿载入。
\section{代数运算：从多张图得到更可靠的图}
\label{sec:algebra-beginner}
本节展开原讲义模块4.4。先不考虑复杂公式：一幅灰度图像就是一张数值表，每个位置存着一个表示明暗的数。所谓\textbf{图像代数运算（image arithmetic）}，就是把两张或多张图\textbf{对应位置上的数}相加、相减、乘系数。例如，同一位置两次读到98和102，把它们平均便得到100。运算对象是数值；能否得到有意义的图像，取决于这些位置是否真的对应同一个物体点。

\subsection{先学会读一张小数值表}
\textbf{矩阵（matrix）}在这里可以先理解为整齐排列的数值网格，并不要求先学线性代数。横着的一排叫\textbf{行（row）}，竖着的一列叫\textbf{列（column）}。本章统一用$r$表示行号，向下增加；用$c$表示列号，向右增加。$f(r,c)$表示图$f$的第$r$行、第$c$列那个数。这样约定是为了避免把“图上的水平边缘”和“数值沿哪个方向变化”混为一谈。
\begin{center}
\begin{tikzpicture}[x=.78cm,y=.78cm,font=\small]
  \draw[step=1] (0,0) grid (3,3);
  \foreach \x/\y/\v in {.5/2.5/20,1.5/2.5/30,2.5/2.5/40,.5/1.5/50,1.5/1.5/60,2.5/1.5/70,.5/.5/80,1.5/.5/90,2.5/.5/100}
    \node at (\x,\y) {\v};
  \draw[very thick] (1,1) rectangle (2,2);
  \draw[-{Stealth}] (0,3.45)--(3,3.45) node[right] {列号$c$增加};
  \draw[-{Stealth}] (-.45,3)--(-.45,0) node[below] {行号$r$增加};
  \node[anchor=west,align=left] at (4,1.7) {若行、列从1开始计数：\\中心像素$f(2,2)=60$\\上邻居$f(1,2)=30$\\右邻居$f(2,3)=70$};
\end{tikzpicture}
\end{center}
本章写出的数值网格既可能代表\textbf{真实像素值}，也可能代表\textbf{待乘的系数}。两者外形相似，含义不同：像素网格告诉我们“图里有什么”，系数网格告诉我们“怎样计算”。

\subsection{逐像素平均：平均的是同一地点的多次观测}
对静止场景连续拍3张照片。假设已对齐，在某个像素位置观测到：
\[
  g_1=96,\qquad g_2=102,\qquad g_3=102.
\]
算术平均（arithmetic mean）就是“先相加，再除以个数”：
\[
  \bar g=\frac{96+102+102}{3}=100.
\]
对每一个位置都做同样的计算，就得到平均图。若共有$N$帧，记号为
\[
  \bar g(r,c)=\frac{1}{N}\sum_{i=1}^{N}g_i(r,c).
\]
这里$\sum$读作“求和”：把$i=1,2,\ldots,N$对应的项全部加起来。上横线只是提示这是平均值，不是另一个坐标。

为什么可能有用？真实场景若应该是100，噪声使一次观测偏小、另一次偏大；平均时这些\textbf{有正有负的偏差有机会相互抵消}。但这不保证任意3张照片都刚好抵消。若3次都是102，平均仍是102。去噪效果是对许多可能的观测而言的统计规律，而不是每次都精确还原。

\textbf{两种平均不要混淆：}
\begin{itemize}
  \item 多帧平均：固定一个空间位置，把\textbf{不同时间}拍到的数平均。对齐准确、场景静止时，不必牺牲空间细节。
  \item 邻域均值滤波：固定一张图，把\textbf{附近不同位置}的数平均。把邻居借来降噪，也可能把不同物体的亮度混在一起；下一节会具体计算。
\end{itemize}

\subsection{噪声、零均值、方差和标准差分别在说什么}
\textbf{噪声（noise）}是观测中不希望出现的扰动。先用最简单的\textbf{加性噪声模型（additive noise model）}：
\[
  g_i=f+n_i.
\]
$f$是真实值，$n_i$是第$i$次误差，$g_i$是拍到的值。真实值100、观测值96，对应$n_i=-4$；观测值103，对应$n_i=3$。为看清推理，这里省略了固定不变的位置$(r,c)$。

\textbf{期望（expectation）}写成$E[\cdot]$，可以先理解为“在同样条件下重复非常多次后，数值趋向的平均”。\textbf{零均值噪声（zero-mean noise）}满足$E[n_i]=0$，意思是长期没有持续偏亮或偏暗的倾向。它\textbf{不表示每一次噪声都是0}，也不表示有限几次噪声的和一定为0。

只知道平均为0，还无法区分“小抖动”和“大抖动”。以下两种噪声长期平均都为0：
\[
  \text{小抖动：}-2,+2,-2,+2,\ldots;
  \quad\text{大抖动：}-10,+10,-10,+10,\ldots
\]
直接平均正负偏差会抵消，所以先把偏差\textbf{平方}，再求平均。这就是\textbf{方差（variance）}：
\[
  \operatorname{Var}(n)=E\bigl[(n-E[n])^2\bigr].
\]
零均值时简化为$E[n^2]$。小抖动的方差是$(4+4)/2=4$，大抖动的方差是$(100+100)/2=100$。\textbf{标准差（standard deviation）}是方差的正平方根：
\[
  \sigma=\sqrt{\operatorname{Var}(n)}.
\]
于是上面两种噪声的标准差分别为2和10。标准差与像素灰度使用同一单位，比较直观；方差使用该单位的平方，推导更方便。\textbf{高斯噪声（Gaussian noise）}指噪声误差的概率分布呈钟形：在零均值模型中，接近0的小误差较常见，较大的正负误差较少见，并且概率关于0对称。它不是说噪声只能取几个固定值，也不是说所有误差概率相同。均值与标准差分别决定钟形的中心位置与宽度。这里讨论已知噪声模型的理论方差，不能把它与“用有限样本估计方差时为何除以$N-1$”混为一件事。

\subsection{为什么平均$N$帧后，标准差只缩小$\sqrt N$倍}
平均图可以拆成“原信号+平均噪声”：
\[
  \bar g=\frac1N\sum_{i=1}^{N}(f+n_i)
        =f+\underbrace{\frac1N\sum_{i=1}^{N}n_i}_{\bar n}.
\]
若各帧都是零均值噪声，那么
\[
  E[\bar g]=f+\frac1N\sum_{i=1}^{N}E[n_i]=f.
\]
这叫\textbf{无偏（unbiased）}：长期平均对准真实值，并非单次结果必然正确。

再看方差。设每帧噪声方差相同，都是$\sigma^2$。对零均值的$\bar n$，先把平方完整展开：
\begin{align*}
 \operatorname{Var}(\bar n)
 &=E\left[\left(\frac{n_1+\cdots+n_N}{N}\right)^2\right]\\
 &=\frac1{N^2}\left(\sum_{i=1}^{N}E[n_i^2]
       +2\sum_{i<j}E[n_i n_j]\right).
\end{align*}
$\sum_{i<j}$表示把不同帧两两配对，每一对只计一次；前面的2来自展开平方时每对会出现$n_in_j$与$n_jn_i$两项。

\textbf{互不相关（uncorrelated）}在这里的关键作用，是让不同帧之间的交叉项$E[n_i n_j]$为0。\textbf{相互独立（independent）}是更强的条件，具有有限方差的独立噪声必然互不相关；本题只需要互不相关。
\[
  \operatorname{Var}(\bar n)
  =\frac{N\sigma^2}{N^2}=\frac{\sigma^2}{N},
  \qquad
  \sigma_{\bar n}=\sqrt{\frac{\sigma^2}{N}}=\frac{\sigma}{\sqrt N}.
\]
\textbf{分母为什么不是$N$？}方差除以$N$以后，标准差还要开平方；把“方差下降”和“标准差下降”说成同一倍率，是这道题最常见的错误。

\begin{center}
\begin{tabular}{rrrr}
\toprule
平均帧数$N$ & 原标准差 & 平均后方差 & 平均后标准差\\
\midrule
1  & 10 & 100 & 10\\
4  & 10 & 25 & 5\\
25 & 10 & 4 & 2\\
100 & 10 & 1 & 1\\
\bottomrule
\end{tabular}
\end{center}
从4帧增加到25帧，成本增加很多，标准差只是从5降到2。因此平均帧数的收益会逐渐变慢。

\subsection{假设不成立时：偏差、相关噪声和运动不会自动消失}
\begin{enumerate}[label=\arabic*.]
  \item \textbf{固定偏差不会被平均掉。}若每次都额外偏亮5，即$E[n_i]=5$，那么$E[\bar g]=f+5$。零均值假设决定了能否无偏。
  \item \textbf{同步波动不能靠重复拍摄抵消。}若各帧都有完全相同的误差$n_i=z$，则$\bar n=z$，方差根本不下降。一般公式中的交叉项此时不能删除。
  \item \textbf{运动会产生重影。}如果第1帧的某位置是路面，第2帧同位置是行人，平均得到的是两种物体的混合；这不是随机噪声抵消。应先\textbf{配准（registration，即对齐对应物体位置）}，或采用适合运动场景的方法。
  \item \textbf{不能先把结果截断再套原推导。}过早四舍五入、裁剪到$0\sim255$会引入额外误差。通常先用更宽的数值类型求和和平均，最后再做显示转换。
\end{enumerate}

\subsection{E3 原题与完整解答：平均去噪需要多少帧}
\label{original:E3}
\textbf{题目。}对同一静止场景连续拍摄$N$帧并逐像素取平均，若各帧噪声互不相关、零均值、方差为$\sigma^2$，求平均后图像的噪声方差与标准差；若要求噪声标准差降为原来的$1/5$，需要平均多少帧？

\textbf{第1步：明确平均对象。}同一位置的真实值$f$保持不变，各帧$g_i=f+n_i$，所以$\bar g=f+\bar n$，其中$\bar n=N^{-1}\sum_i n_i$。

\textbf{第2步：代入题设。}零均值给出$E[\bar n]=0$；互不相关使交叉项消失。因此
\[
  \boxed{\operatorname{Var}(\bar n)=\frac{\sigma^2}{N}},
  \qquad \boxed{\sigma_{\bar n}=\frac{\sigma}{\sqrt N}}.
\]
\textbf{第3步：把文字要求变成等式。}要求的是\emph{标准差}降到$1/5$，不是方差降到$1/5$：
\[
 \frac{\sigma}{\sqrt N}=\frac{\sigma}{5}
 \ \Longrightarrow\ \sqrt N=5
 \ \Longrightarrow\ \boxed{N=25\text{帧}}.
\]
若要求“至少下降到原来的$1/5$”，应写$N\ge25$，最少25帧。例如原标准差10、原方差100，平均25帧后标准差2、方差4；$2/10=1/5$而$4/100=1/25$，互相核对一致。

\subsection{图像相减：差别会出现，负数也会出现}
\textbf{图像相减（image subtraction）}逐位置计算$d=f-h$。例如一个位置前后分别为100和130，按“后减前”得到30；另一处由100变成70，差值为$-30$。正负号携带“变亮或变暗”的信息。

它可用于比较前后变化、背景差分以及后面的细节提取，但要注意三件事：
\begin{itemize}
  \item 两图应尺寸一致、位置对齐，且曝光变化要受控；否则整幅图的差异未必来自物体运动。
  \item 用可容纳负数的类型保留中间结果。无符号8位数只能存$0\sim255$，直接相减可能发生裁剪或回绕，具体取决于软件。
  \item 想显示差值，可以取绝对值$|d|$观察变化幅度，或把0显示为中灰以区分正负；这都是\textbf{显示规则}，不等于原始差值计算。
\end{itemize}
相减也会增加不相关噪声：两图噪声方差若都是$\sigma^2$，则差值噪声方差为$\sigma^2+\sigma^2=2\sigma^2$。减法强调差异，却不天然具有去噪能力。

\section{空间滤波：让邻居一起决定一个像素}
\label{sec:spatial-filter-beginner}
本节展开原讲义模块4.5。\textbf{空间域（spatial domain）}就是直接在像素网格上操作。\textbf{滤波（filtering）}不是只指“去掉噪声”，也包括平滑、边缘提取和锐化。共同点是：为一个输出位置收集附近信息，按指定规则得到一个新数值。

\subsection{邻域、窗口、核：三个名字对应什么}
\textbf{邻域（neighborhood）}是一组附近像素；\textbf{窗口（window）}规定选取它们的位置与形状。常见$3\times3$窗口含中心及周围8个像素。$3\times3$中间有唯一中心；这就是教学中常用奇数边长的原因。

\textbf{核（kernel）}也叫\textbf{掩模（mask）}或模板。在\textbf{线性滤波（linear filtering）}中，它是一张与窗口同形状的系数表。把系数表覆盖到像素表上，\textbf{对应位置相乘，再把所有乘积相加}，得到一个输出像素。这里不是线性代数里的矩阵乘法：不需要“行乘列”，只需要“同位置乘后相加”。
\begin{center}
\begin{tikzpicture}[x=.62cm,y=.62cm,font=\small]
  \foreach \x/\y/\v in {0/2/10,1/2/20,2/2/30,0/1/40,1/1/50,2/1/60,0/0/70,1/0/80,2/0/90}
   {\draw (\x,\y) rectangle ++(1,1);\node at (\x+.5,\y+.5) {\v};}
  \node at (1.5,3.65) {像素值$P$};
  \node at (4,1.5) {$\longrightarrow$};
  \foreach \x/\y/\v in {5/2/1,6/2/2,7/2/1,5/1/2,6/1/4,7/1/2,5/0/1,6/0/2,7/0/1}
   {\draw (\x,\y) rectangle ++(1,1);\node at (\x+.5,\y+.5) {\v};}
  \node[align=center] at (6.5,3.65) {对应相乘的系数\\最后整体除以16};
  \draw[-{Stealth}] (8.5,1.5)--(10,1.5);
  \draw (10.4,1) rectangle (11.4,2);
  \node at (10.9,1.5) {50};
  \node[align=center] at (10.9,-.45) {9个输入数\\得到1个输出数};
\end{tikzpicture}
\end{center}
以上计算全部展开为
\begin{align*}
 g&=\frac1{16}(1\cdot10+2\cdot20+1\cdot30
       +2\cdot40+4\cdot50+2\cdot60\\
  &\hspace{23mm}+1\cdot70+2\cdot80+1\cdot90)\\
  &=\frac{10+40+30+80+200+120+70+160+90}{16}
   =\frac{800}{16}=50.
\end{align*}
中心50得到较大的系数4，左右上下分别得到系数2，四个角得到系数1。\textbf{系数的大小表示这个邻居的影响有多大}；负系数表示从总和中减去该邻居的贡献。

\subsection{滑动一次只写一个数：输入图和输出图要分开}
算出一个中心的结果以后，把窗口移动到下一个中心，再取新的邻居重复计算。下图中粗框表示两个相邻窗口的位置，并非一次把整块$3\times3$都改成同一个数。
\begin{center}
\begin{tikzpicture}[x=.65cm,y=.65cm,font=\small]
 \draw[step=1,gray] (0,0) grid (6,4);
 \foreach \x in {0,...,5} \foreach \y in {0,...,3}
  {\pgfmathtruncatemacro{\v}{10+\x*5+\y*3}\node at (\x+.5,\y+.5) {\v};}
 \draw[very thick] (0,1) rectangle (3,4);
 \draw[thick,dashed] (1,1) rectangle (4,4);
 \draw[very thick] (1,2) rectangle (2,3);
 \draw[thick,dashed] (2,2) rectangle (3,3);
 \draw[-{Stealth},thick] (1.5,4.5)--(2.5,4.5);
 \node[align=left,anchor=west] at (7,2.1) {实线：先计算这个中心\\虚线：再向右移动一格\\每次读取原输入图\\各结果写到另一张输出图};
\end{tikzpicture}
\end{center}
如果刚算完就覆盖原图，再让后续窗口读到这个新值，后续计算便混合了原值与已修改值；从左向右算和从右向左算甚至可能得到不同结果。普通空间滤波默认\textbf{所有窗口都读取同一份未被本次运算改写的输入}。

操作步骤可以直接照着执行：
\begin{enumerate}[label=\arabic*.]
 \item 准备输出图，并选定窗口大小、核和边界规则。
 \item 把核的中心对准一个输入像素，取出所需邻域。
 \item 对应相乘，再相加；若采用中值滤波，则改为排序取中间数。
 \item 把结果写在输出图的同一中心位置，移动窗口，直到遍历完成。
 \item 最后才按要求取整、裁剪或转换为显示格式。
\end{enumerate}

\subsection{相关与卷积：到底要不要翻转系数表}
\textbf{相关（correlation）}按原样使用核。设核$w(s,t)$中，$s$表示相对中心的行偏移，$t$表示列偏移：负1是上一行或左一列，0是本行或本列，正1是下一行或右一列。对$3\times3$核，
\[
 g(r,c)=\sum_{s=-1}^{1}\sum_{t=-1}^{1}w(s,t)f(r+s,c+t).
\]
两个求和号只是“逐行、逐列把9项加起来”。

\textbf{卷积（convolution）}则在上述乘加之前把核转$180^\circ$，也就是上下翻转并左右翻转：
\[
 g(r,c)=\sum_{s=-1}^{1}\sum_{t=-1}^{1}w(s,t)f(r-s,c-t).
\]
等价做法是把$w(s,t)$换成$w(-s,-t)$，再按相关的方法计算。\textbf{只上下翻转或只左右翻转都不是一般二维卷积所需的$180^\circ$翻转。}

例如取一行像素$[10,20,40]$，核$[1,0,-1]$：
\[
 \begin{array}{ll}
 \text{相关：}&1\cdot10+0\cdot20-1\cdot40=-30,\\
 \text{卷积：}&-1\cdot10+0\cdot20+1\cdot40=30.
 \end{array}
\]
同样的边界变化，输出符号相反。对于$180^\circ$旋转对称的均值核、高斯核、拉普拉斯核，翻转后不变，因此相关与卷积给出相同结果。本章A1、A2的7个核都具有这种对称性，题内计算不受这一区别影响；以后遇到一阶差分核时必须检查。

\subsection{走到图像边缘怎么办：延拓也是算法的一部分}
\textbf{边界延拓／填充（padding）}是为图外不存在的像素约定临时取值。比如图像最左侧像素要找“左邻居”，总得规定这个值是什么。下面用一行$[10,20,30,40]$说明，两边各补2个数，竖线中间才是真实数据。
\begin{center}
\begin{tabular}{ll}
\toprule
规则 & 延拓后的数列\\
\midrule
常数填充（constant，取0） & $0,0\mid10,20,30,40\mid0,0$\\
复制边缘（replicate） & $10,10\mid10,20,30,40\mid40,40$\\
对称反射，重复端点 & $20,10\mid10,20,30,40\mid40,30$\\
反射，不重复端点 & $30,20\mid10,20,30,40\mid30,20$\\
周期延拓（circular） & $30,40\mid10,20,30,40\mid10,20$\\
\bottomrule
\end{tabular}
\end{center}
不同软件对“symmetric”和“reflect”的命名细节可能不同，所以\textbf{看到名字后还要核对是否重复端点}。二维情况按行、列同样延伸，四角也要服从同一规则。

用3点均值计算最左边那个10，可以立刻看出差异：
\[
 \begin{array}{ll}
 \text{补0：}&(0+10+20)/3=10,\\
 \text{复制：}&(10+10+20)/3=40/3\approx13.33,\\
 \text{不重复端点的反射：}&(20+10+20)/3=50/3\approx16.67,\\
 \text{周期：}&(40+10+20)/3=70/3\approx23.33.
 \end{array}
\]
核完全一样，边缘输出仍可不同。周期延拓会把最右端接到最左端，适合本来就周期重复的内容；对普通照片，它可能创造不真实的边界跳变。补0会在亮图周围引入暗边，因此“系数和为1保持常数”在图像边缘还需要\textbf{延拓本身也保持这个常数}。

另一种选择是\textbf{只计算窗口完全落在图内的位置（valid输出）}。例如$5\times5$图使用$3\times3$核、不填充时，只能得到$3\times3$个完整输出。若要输出与输入一样大，通常需要某种填充。

\subsection{先做常数图测试：系数和告诉我们什么}
把窗口内所有像素都设成同一个数$C$，如100。每个数都一样，因此无论窗口如何移动，输出都是
\[
 g=Cw_1+Cw_2+\cdots+Cw_K
   =C\underbrace{(w_1+\cdots+w_K)}_{S},
\]
其中$K$是系数个数，$S$是系数和。这个最简单的测试能回答“均匀背景会发生什么”：
\begin{itemize}
 \item $S=1$：100仍输出100，\textbf{保留常数亮度}。
 \item $S=0$：100输出0，\textbf{去除常数背景}，只对某些变化产生响应。
 \item $S=2$：100输出200，常数背景被放大2倍。
\end{itemize}
以后频域中把不随位置变化的成分叫\textbf{直流分量（DC component）}；这里无需先懂傅里叶变换，就能理解系数和是直流增益。

\textbf{重要边界：系数和不是万能分类器。}和为1不一定是平滑核，也可能是“原图+细节”的最终锐化核；和为0只证明抑制常数，不能仅凭这一条就保证所有高频都通过。还应看正负系数的布局、对图案的实际响应，必要时再研究完整频率响应。

\subsection{均值滤波：每个邻居同样重要}
\textbf{均值滤波（mean filter）}对窗口里所有数赋予相等权重。$3\times3$中有9个数，核为
\[
 B=\frac19\begin{bmatrix}1&1&1\\1&1&1\\1&1&1\end{bmatrix}.
\]
它的系数和是$9/9=1$。假设小区域本来处处是50，观测为
\[
 P=\begin{bmatrix}48&52&49\\51&50&53\\47&51&49\end{bmatrix},
 \qquad
 g=\frac{48+52+49+51+50+53+47+51+49}{9}=50.
\]
这次正负噪声刚好抵消；其他观测未必精确得到50。在局部真实亮度相同、9个噪声零均值且互不相关、方差相同的条件下，输出噪声方差为$\sigma^2/9$，标准差为$\sigma/3$。推理与多帧平均一样，但\textbf{现在假设的是空间邻居之间的噪声关系和真实亮度关系}。

为什么会模糊边缘？如果左边物体灰度0、右边物体灰度90，窗口为
\[
 P=\begin{bmatrix}0&0&90\\0&0&90\\0&0&90\end{bmatrix},
 \qquad g=(0+0+90+0+0+90+0+0+90)/9=30.
\]
中心原本是0，却被右边物体拉到30。\textbf{平均不认识物体边界，只认识窗口内的数。}窗口越大，随机噪声通常更易平滑，但越容易混合相隔更远的区域。

\subsection{高斯滤波：中心近的像素多说一点话}
\textbf{高斯滤波（Gaussian filter）}仍然做乘加，但采用中心大、远处小的权重。常用的小型近似核为
\[
 B_G=\frac1{16}\begin{bmatrix}1&2&1\\2&4&2\\1&2&1\end{bmatrix}.
\]
其权重总和$1+2+1+2+4+2+1+2+1=16$，除以16后和为1。对称布局没有特别偏向左边或右边；中心权重$4/16$比角点$1/16$大。

“高斯”这个名字来自钟形权重曲线：离中心越远，权重逐渐下降。更大的离散核常由
\[
 \widetilde w(s,t)=\exp\!\left(-\frac{s^2+t^2}{2\sigma_G^2}\right),
 \qquad
 w(s,t)=\frac{\widetilde w(s,t)}{\sum_{s,t}\widetilde w(s,t)}
\]
构造。$s^2+t^2$是到中心的距离平方；$\sigma_G>0$控制钟形有多宽；若$s,t$以像素间隔计，它也以像素间隔为单位。\textbf{它是滤波器参数，不是前面噪声的标准差$\sigma$}。第二步归一化保证有限窗口内系数和为1。

对平坦区域中互不相关、同方差的噪声，一般加权平均的噪声方差为
\[
 \operatorname{Var}\!\left(\sum_j w_jn_j\right)=\sigma^2\sum_jw_j^2.
\]
上面的$3\times3$高斯近似核给出
\[
 \sigma^2\frac{1^2+2^2+1^2+2^2+4^2+2^2+1^2+2^2+1^2}{16^2}
 =\frac{36}{256}\sigma^2=\frac9{64}\sigma^2.
\]
同样9个独立样本，等权均值的方差$\sigma^2/9$更小；高斯权重选择的是\textbf{更重视近邻的空间折中}，不能声称它在所有条件下都比均值去噪更强。两者都会模糊边缘；“高斯”也不意味着任何高斯噪声都能无损去除。

\subsection{中值滤波：先排队，再取中间那个数}
\textbf{中值／中位数（median）}不是平均数。把数从小到大排列，取最中间的一个。9个数排序后取第5个；$5\times5$窗口有25个数，就取第13个。

下面是一块真实参与计算的9数窗口，中心遭遇了很亮的脉冲：
\[
 P=\begin{bmatrix}10&12&11\\9&255&10\\11&8&10\end{bmatrix}.
\]
\textbf{均值的算法：}把9个数相加，得到
\[
 10+12+11+9+255+10+11+8+10=336,
 \qquad\text{均值}=336/9\approx37.33.
\]
\textbf{中值的算法：}先把它们抄成一列并排序：
\[
 8,\ 9,\ 10,\ 10,\ \underbrace{10}_{\text{第5个}},\ 11,\ 11,\ 12,\ 255,
 \qquad\text{中值}=10.
\]
输出只把\textbf{中心位置}改成10，并不是把整个$3\times3$都改成10。255很大，却位于排序队伍末端，未能主导中值；均值则被它明显拉高。

若同一块窗口使用高斯近似核，中心脉冲恰好得到最大权重：
\[
 \frac{10+24+11+18+1020+20+11+16+10}{16}
 =\frac{1140}{16}=71.25.
\]
这说明面对孤立极值，简单加权平均不一定是合适选择。\textbf{椒盐噪声（salt-and-pepper noise）}常表现为随机黑点与白点；更一般的\textbf{脉冲噪声（impulse noise）}指少数位置突然出现很大的异常值。中值滤波常适合处理这类异常。

\subsection{中值为什么常保边，却也会删掉细线}
比较两种$3\times3$窗口，灰度只取0与90：
\[
 P_{\text{边缘}}=\begin{bmatrix}0&0&90\\0&0&90\\0&0&90\end{bmatrix},
 \qquad
 P_{\text{细线}}=\begin{bmatrix}0&90&0\\0&90&0\\0&90&0\end{bmatrix}.
\]
第一块中6个0、3个90，中值为0，中心暗像素保持0；均值则为30，产生过渡灰。第二块也是6个0、3个90，\textbf{中值同样为0}，但此时中心原本属于真实亮线，亮线被删除了。算法仅凭窗口数值无法自动知道哪三个90是细线、哪三个90是噪声。

\begin{center}
\begin{tikzpicture}[x=.55cm,y=.55cm,font=\small]
 \foreach \x in {0,...,5} \foreach \y in {0,...,3}{
   \ifnum\x<3 \fill[white] (\x,\y) rectangle ++(1,1);
   \else \fill[gray!55] (\x,\y) rectangle ++(1,1);\fi}
 \fill[black] (1,2) rectangle (2,3);\fill[white] (4,1) rectangle (5,2);
 \draw[step=1] (0,0) grid (6,4);
 \node at (3,-.8) {带脉冲的边界};
 \draw[-{Stealth},thick] (6.7,2)--(8,2);
 \foreach \x in {9,...,14} \foreach \y in {0,...,3}{
   \ifnum\x<12 \fill[white] (\x,\y) rectangle ++(1,1);
   \else \fill[gray!55] (\x,\y) rectangle ++(1,1);\fi}
 \draw[step=1] (9,0) grid (15,4);
 \node at (12,-.8) {去除孤立异常示意};
 \node[align=center,text width=12cm] at (7.5,-2.5) {是否保留结构取决于窗口内的数量与形状\\窄线、小点、尖角都可能被中值改变};
\end{tikzpicture}
\end{center}
所以“中值保边”是相对于某些线性平均方法的常见优势，\textbf{并非保护任意边缘、细线或角点的保证}。噪声比例很大、窗口过大、目标本身太小时，结果都可能变差。

\subsection{最大值与最小值：排序后也可以选两端}
\textbf{顺序统计滤波（order-statistic filtering）}先根据数值大小排列窗口，再选择某个次序的值。中值选择中间；\textbf{最大值滤波（maximum filter）}选择最后一个，\textbf{最小值滤波（minimum filter）}选择第一个。

仍用窗口$[10,12,11,9,255,10,11,8,10]$，排序后最小值为8、中值为10、最大值为255。这三种操作得到不同答案，因为它们在回答不同的问题：窗口中“最暗是多少”“中间是多少”“最亮是多少”。最大值会让亮区域向邻近位置扩展，有机会填掉孤立暗点，却会扩大亮脉冲；最小值反过来，让暗区域扩展，可去掉孤立亮点，却会扩大暗脉冲。因此二者并不是与中值等价的通用去椒盐方法。

在灰度形态学中，采用同样的\textbf{平坦结构元素（flat structuring element，即只指定参与比较的位置、不另加高度权重）}时，局部最大值对应\textbf{膨胀（dilation）}，局部最小值对应\textbf{腐蚀（erosion）}，具体位置还需遵守所采用的反射／锚点约定。本章只需记住数值规则与亮暗区域的变化，不必先学习完整形态学。

\subsection{中值为什么不是固定系数卷积}
线性方法应当满足“先相加再滤波”等于“分别滤波再相加”。用3点中值就能构造反例：
\[
 a=[0,0,1],\quad b=[0,1,0],\quad a+b=[0,1,1].
\]
分别取中值：$\operatorname{median}(a)=0$，$\operatorname{median}(b)=0$，两者之和为0；先相加再取中值却得到1。因此
\[
 \operatorname{median}(a+b)\ne\operatorname{median}(a)+\operatorname{median}(b).
\]
这就是\textbf{非线性（nonlinear）}的具体证据。中值可以用滑动窗口规定“看哪些像素”，但不能用一张对所有输入都固定不变的系数表实现同样操作。

\subsection{先看变化量，再认识导数}
一排像素为$[0,0,0,10,10,10]$，从左往右在第3和第4个像素之间由暗变亮。相邻两数相减，得到
\[
 [0-0,\ 0-0,\ 10-0,\ 10-10,\ 10-10]=[0,0,10,0,0].
\]
这叫\textbf{一阶差分（first difference）}：问的是“下一点比这一点变了多少”。平坦部分给0，跳变位置给10。

再比较相邻的变化量，便得到\textbf{二阶差分（second difference）}。一个常用的中心写法为
\[
 L_j=f_{j-1}-2f_j+f_{j+1}.
\]
注意中心要乘$-2$，不是只把两个邻居相加。它也可以理解为
\[
 L_j=(f_{j+1}-f_j)-(f_j-f_{j-1}),
\]
即“右边的变化量减左边的变化量”。在上面的阶跃两侧，中心取第3个像素时$L=0-2\cdot0+10=10$；中心取第4个像素时$L=0-2\cdot10+10=-10$。于是出现一正一负的成对响应。
\begin{center}
\begin{tikzpicture}[x=.83cm,y=.055cm,font=\small]
 \draw[-{Stealth}] (0,-5)--(6.7,-5) node[right] {位置};
 \draw[thick] (.5,0)--(2.5,0)--(3.5,10)--(5.5,10);
 \foreach \x/\v in {.5/0,1.5/0,2.5/0,3.5/10,4.5/10,5.5/10}{\fill (\x,\v) circle (1.7pt);}
 \node[anchor=east] at (-.1,8) {原像素};
 \draw[gray] (0,-38)--(6,-38);
 \draw[thick] (3,-38)--(3,-28);\fill (3,-28) circle (1.7pt);
 \node[anchor=east] at (-.1,-32) {一阶：$+10$};
 \draw[gray] (0,-78)--(6,-78);
 \draw[thick] (2.5,-78)--(2.5,-68);\fill (2.5,-68) circle (1.7pt);
 \draw[thick] (3.5,-78)--(3.5,-88);\fill (3.5,-88) circle (1.7pt);
 \node[anchor=east] at (-.1,-77) {二阶：$+10,-10$};
 \node[align=left,anchor=west] at (7,-40) {一阶定位变化\\二阶定位变化量的改变\\正负号不能随意丢掉};
\end{tikzpicture}
\end{center}
再看匀速上升的数列$[0,10,20,30,40]$。一阶差分都是10；中间各点的二阶差分都是0，例如$10-2\cdot20+30=0$。\textbf{二阶并非“比一阶更大的同一个数”，而是在问另一件事。}

连续数学中，一阶变化率叫\textbf{导数（derivative）}；图像里沿行、列方向分别计算的一阶变化量组成\textbf{梯度（gradient）}。像素网格上的差分是导数的离散近似。本章先掌握上面的减法，就能读懂后续导数符号。

\subsection{方向要说完整：沿列变化与竖直走向的边缘}
考虑两块窗口：
\[
 P_c=\begin{bmatrix}10&10&90\\10&10&90\\10&10&90\end{bmatrix},
 \qquad
 P_r=\begin{bmatrix}10&10&10\\10&10&10\\90&90&90\end{bmatrix}.
\]
$P_c$从左到右，也就是\textbf{沿列号$c$增加的方向}变化；边界本身却是一条\textbf{竖直走向}的线。$P_r$从上到下，也就是\textbf{沿行号$r$增加的方向}变化；边界本身是\textbf{水平走向}的线。

例如核
\[
 K_c=\begin{bmatrix}-1&0&1\\-1&0&1\\-1&0&1\end{bmatrix}
\]
按相关计算，对$P_c$得到$3(90-10)=240$，对$P_r$得到0。这类核比较右列和左列，因此对\textbf{左右亮度变化}敏感，常用来发现\textbf{竖直走向边界}。只写“水平算子”或“垂直边缘算子”容易因命名习惯不同而歧义，答题时应写清楚比较哪两侧。

\subsection{拉普拉斯：把两个方向的二阶差分加起来}
\textbf{拉普拉斯算子（Laplacian operator）}把行方向、列方向的二阶差分相加。这里默认相邻采样间隔为1；若要估计真实物理坐标中的二阶导数，还应分别除以对应采样间隔的平方。先固定一种符号约定：
\begin{align*}
 Lf(r,c)={}&f(r-1,c)+f(r+1,c)+f(r,c-1)+f(r,c+1)\\
          &-4f(r,c).
\end{align*}
这一写法近似连续的$\nabla^2 f$，对应中心为\textbf{负4}的核：
\[
 L=\begin{bmatrix}0&1&0\\1&-4&1\\0&1&0\end{bmatrix}.
\]
这里只使用上、下、左、右4个邻居，叫\textbf{四邻域（4-neighborhood）}。若定义相反符号$P=-L$，核就变成
\[
 P=\begin{bmatrix}0&-1&0\\-1&4&-1\\0&-1&0\end{bmatrix}.
\]
两种写法都可以，但\textbf{公式、核与最后加减号必须配套}。

设中心70、四邻居都是50：
\[
 Lf=50+50+50+50-4\cdot70=-80,\qquad Pf=80.
\]
中心比周围亮，$Lf$为负；使用$g=f-Lf$时，中心变成$70-(-80)=150$，比原来更亮，局部差异增强。同一个结果也可以写成$g=f+Pf=70+80=150$。
\[
 \boxed{g=f-Lf=f+Pf},\qquad
 \boxed{K_{\rm sharp}=I-L=I+P=
 \begin{bmatrix}0&-1&0\\-1&5&-1\\0&-1&0\end{bmatrix}}.
\]
$I$不是把所有数都设成1的网格，而是\textbf{只复制中心}的单位脉冲核：
\[
 I=\begin{bmatrix}0&0&0\\0&1&0\\0&0&0\end{bmatrix}.
\]
纯拉普拉斯核的系数和为0，常数区消失；最终锐化核的系数和为1，常数区保留。\textbf{“提取细节图”和“输出锐化照片”是两个不同阶段。}

加入对角邻居的常见正中心版本为
\[
 P_8=\begin{bmatrix}-1&-1&-1\\-1&8&-1\\-1&-1&-1\end{bmatrix}.
\]
它对8个方向对称取样，常称八邻域拉普拉斯类核。四邻域和八邻域版本都没有选定单一朝向，但在有限方格上\textbf{不具有连续平面任意角度下的严格旋转不变性}。考试可说“近似无单一方向偏好”，不要把它理解成各方向的数值响应必然完全一致。

\subsection{一阶、二阶与噪声：不是越尖锐越好}
灰度阶跃在一阶差分中常出现单个主要响应，在二阶差分中常出现一正一负的双响应。二阶对细线、孤立点等快速局部变化很敏感，适合构造细节增强；但孤立噪声也是快速变化，因此也会被放大。

“梯度边缘更粗、二阶边缘更细”只是对某些离散算子与显示处理的常见比较，不是所有核、尺度、阈值下都成立的定理。例如上一幅理想阶跃的一阶前向差分只有一个非零样本，二阶反而有两个。答题时应把\textbf{阶跃响应、细节敏感性和噪声敏感性}说清，不要仅背“粗”和“细”。

最终锐化也可能产生\textbf{过冲／欠冲（overshoot/undershoot）}。对一行$[0,0,0,10,10,10]$使用$g=f-Lf$，阶跃附近可出现$-10$和20，超出原范围$[0,10]$。如果直接裁剪，会改变本来推导的线性结果；如果增益很大，会在边界两侧看见不自然的明暗边。实际处理中应控制锐化强度，并考虑先抑制噪声。

\subsection{非锐化掩模：先模糊，再相减，最后加回}
\textbf{非锐化掩模（unsharp masking）}这个名字容易误导：目的是让图变清晰，但需要先制作一张“不锐”的模糊图。设原图为$f$，用归一化平滑核得到$f_{\rm blur}$，则
\[
 m=f-f_{\rm blur}
\]
是\textbf{细节掩模（detail mask）}。平坦区域中两者近似相同，相减接近0；细节处不同，相减留下正负细节。接着把细节按强度$k$加回原图：
\[
 \boxed{g=f+k\,m=f+k(f-f_{\rm blur})},\qquad k\ge0.
\]
$k=0$不改变图像，$k=1$是常见的基本锐化，$k$越大，细节放大通常越明显，噪声与明暗边也可能越明显。
\begin{center}
\begin{tikzpicture}[node distance=9mm and 13mm,font=\small,
 box/.style={draw,rounded corners=1mm,minimum height=9mm,align=center,inner sep=4pt}]
 \node[box] (f) {原图$f$};
 \node[box,right=of f] (blur) {平滑\\$f_{\rm blur}$};
 \node[box,right=of blur] (sub) {相减\\$m=f-f_{\rm blur}$};
 \node[box,below=of sub] (gain) {乘强度$k$};
 \node[box,left=of gain] (add) {加回原图\\$g=f+km$};
 \draw[-{Stealth}] (f)--(blur);
 \draw[-{Stealth}] (blur)--(sub);
 \draw[-{Stealth}] (f.north)--++(0,.7)-|(sub.north);
 \draw[-{Stealth}] (sub)--(gain);
 \draw[-{Stealth}] (gain)--(add);
 \draw[-{Stealth}] (f.south)|-(add.west);
\end{tikzpicture}
\end{center}
数值例子：一行灰度为$[40,40,40,80,80,80]$，用3点均值作模糊。在暗侧紧靠边缘的40处，
\[
 f_{\rm blur}=(40+40+80)/3=53.33\ldots,\quad
 m=40-53.33\ldots=-13.33\ldots;
\]
在亮侧紧靠边缘的80处，
\[
 f_{\rm blur}=(40+80+80)/3=66.66\ldots,\quad
 m=80-66.66\ldots=13.33\ldots.
\]
取$k=1$，两个边缘像素变成约26.67和93.33，明暗差加大；离边缘较远处仍为40或80。这说明锐化改变的是\textbf{局部对比}，不是凭空恢复相机没拍到的真实细节。

\subsection{高提升的两种写法：必须先检查亮度是否归一化}
\textbf{高提升滤波（high-boost filtering）}的记号在教材间并不完全一致。一种常用写法是上述$g=f+k(f-f_{\rm blur})$，较大的$k$表示更强的细节提升。原讲义使用的是另一种形式：
\[
 g_A=A f-f_{\rm blur}=(A-1)f+(f-f_{\rm blur}),\qquad A>1.
\]
它们都涉及原图和模糊图，但参数\textbf{不能不经换算就当作同一个量}。用常数图$f=C$测试，归一化模糊给出$f_{\rm blur}=C$，于是
\[
 f+k(f-f_{\rm blur})=C,\qquad g_A=(A-1)C.
\]
例如均匀区$C=80$：
\[
 A=1.5\Rightarrow g_A=40;\quad
 A=2\Rightarrow g_A=80;\quad
 A=3\Rightarrow g_A=160.
\]
所以\textbf{$Af-f_{\rm blur}$并非对任意$A>1$都保持亮度}；它只有$A=2$时保持常数，除非另做归一化。在$A=1$时，得到的只是$m=f-f_{\rm blur}$，不是已把细节加回原图的最终锐化图。

若将$g_A$除以$A-1$，则
\[
 \frac{g_A}{A-1}
 =f+\frac1{A-1}(f-f_{\rm blur}),\qquad A>1.
\]
它才具有常数增益1，相当于$k=1/(A-1)$。也就是说，\textbf{按这种归一化约定，$A$越大，细节相对原图的提升反而越小}。背公式时要同时说出输出是否归一化，不能只记“参数越大越锐”。

\subsection{A1 原题：四个掩模逐一计算与判读}
\label{original:A1}
\textbf{题目【练习1原题】。}指出下列掩模作用于图像，会产生什么特性。
\[
 \begin{array}{ll@{\qquad}ll}
 \mathrm{a.}&\begin{bmatrix}2&-1&-1\\-1&2&-1\\-1&-1&2\end{bmatrix}&
 \mathrm{b.}&\begin{bmatrix}0&-1&0\\-1&4&-1\\0&-1&0\end{bmatrix}\\[1.1em]
 \mathrm{c.}&\dfrac12\begin{bmatrix}-2&1&-2\\1&6&1\\-2&1&-2\end{bmatrix}&
 \mathrm{d.}&\dfrac19\begin{bmatrix}1&1&1\\1&1&1\\1&1&1\end{bmatrix}
 \end{array}
\]
判读顺序是：\textbf{先算系数和，再看系数布局，再用简单图案验证}。外面的$1/2$或$1/9$必须乘到每个系数上。

\subsubsection{A1(a)：方向选择性的零和核，不能只说“对角线”}
系数和为$3\cdot2+6\cdot(-1)=0$，所以常数背景输出0。更特别的是：\textbf{每一行、每一列的系数和也都为0}。

如果窗口每一行内部都是相同值，例如
\[
 P=\begin{bmatrix}10&10&10\\20&20&20\\40&40&40\end{bmatrix},
\]
输出为$10(2-1-1)+20(-1+2-1)+40(-1-1+2)=0$。同理，若每一列内部都相同，也输出0。因此它会抑制只沿行方向或只沿列方向变化的这种简单图案。

再对比两条斜亮线：
\[
 P_{\backslash}=\begin{bmatrix}10&0&0\\0&10&0\\0&0&10\end{bmatrix},
 \qquad
 P_{/}=\begin{bmatrix}0&0&10\\0&10&0\\10&0&0\end{bmatrix}.
\]
\[
 \begin{array}{ll}
 K_a\text{作用于}P_{\backslash}:&2\cdot10+2\cdot10+2\cdot10=60,\\
 K_a\text{作用于}P_{/}:&-1\cdot10+2\cdot10-1\cdot10=0.
 \end{array}
\]
这说明同一窗口中心处，左上到右下亮线和右上到左下亮线的响应不同。\textbf{这只是给定图案、给定对齐位置的响应，不能推出另一方向的整条线在所有输出位置都为0。}

\textbf{结论：}这是有方向选择性的零和差分／高通类核，突出某些斜向细节，抑制常数及上述单方向变化。它的$180^\circ$对称布局不同于反对称的一阶差分；若用平滑函数的差分近似分析，其低阶主项属于混合二阶变化。零基础答题时，用“对哪些明确排列响应强、对哪些排列为0”比空泛写“增强对角线”更准确。

\subsubsection{A1(b)：四邻域正中心拉普拉斯核}
系数和$4-1-1-1-1=0$。它计算
\[
 g=4f_{\rm 中}-f_{\rm 上}-f_{\rm 下}-f_{\rm 左}-f_{\rm 右}.
\]
常数50输出$4\cdot50-4\cdot50=0$。若中心70、四邻居50，则输出$4\cdot70-4\cdot50=80$。若中心30、四邻居50，则输出$120-200=-80$。正负响应对应中心相对邻居的亮暗差异。

\textbf{结论：}它是$P=-L$，用于提取二阶细节、细线、孤立点及边界变化，均匀区归零。没有单一选定朝向，但离散方向响应并非严格各向相同。若要得到保留背景的锐化图，应再加原图$g=f+Pf$，不能把这个零和输出直接叫作保亮度的最终照片。

\subsubsection{A1(c)：和为1的细节强调核，不是均值平滑}
先把所有整数相加：4个角贡献$-8$，上下左右贡献4，中心贡献6，总和2。再乘外部$1/2$：
\[
 S=\frac12(-8+4+6)=1.
\]
所以常数50仍输出50。中心实际系数为3，四角为$-1$，轴向邻居为$1/2$，具有明显的减法与中心强调；不能因为和为1就判为低通。

对中心70、其余8个邻居50的窗口：
\[
 g=\frac12\left[6\cdot70+(50+50+50+50)
                     -2(50+50+50+50)\right]
  =\frac{420+200-400}{2}=110.
\]
中心从70变成110，平坦背景50保持50。也可以把它分成$I+(K_c-I)$，其中后一项系数和为0，表示“原图+某种细节”。

\textbf{结论：}这是保留常数背景的锐化／细节强调类核，不是纯高通，也不是等权平均。其轴向邻居为正、对角邻居为负，具体频率与方向响应由完整布局决定；不能说“每一种高频都同样增强”，也不能不加验证地说它就是某个非负模糊核的标准高提升形式。

\subsubsection{A1(d)：$3\times3$均值低通核}
系数和$9\cdot(1/9)=1$，所有系数相等且非负。每个输出都是邻域9个数的平均。常数50仍为50；前文8个近10的数加一个255的窗口，则输出$336/9\approx37.33$。

\textbf{结论：}这是均值平滑／低通滤波器，可减弱局部随机起伏，也会模糊边缘和小结构。窗口跨过0到90的边界时会产生30、60等中间灰度；对强脉冲噪声不如合适的中值处理稳健。

\subsection{A2 原题：三种核的性质与方向}
\label{original:A2}
\textbf{题目。}判断下列三个掩模各属于低通还是高通、平滑还是锐化、是否有方向性。
\[
 \text{① }\begin{bmatrix}0&-1&0\\-1&5&-1\\0&-1&0\end{bmatrix}
 \qquad
 \text{② }\begin{bmatrix}-1&-1&-1\\-1&8&-1\\-1&-1&-1\end{bmatrix}
 \qquad
 \text{③ }\begin{bmatrix}-1&-1&-1\\2&2&2\\-1&-1&-1\end{bmatrix}.
\]

\subsubsection{A2①：原图加四邻域细节，最终锐化核}
$S=5-4=1$，故常数50输出50。它恰好等于$I+P$，所以是\textbf{保留常数背景的高频强调／锐化核}，不是去掉背景的纯高通输出。

中心70、上下左右50时：
\[
 g=5\cdot70-50-50-50-50=350-200=150.
\]
没有选定唯一方向，使用四邻域的近似方向对称结构。但如果题目要求二选一“低通或高通”，应说明它属于\textbf{锐化的高频强调类，同时保留直流}，而不是因为系数和1就选低通。

\subsubsection{A2②：八邻域细节提取核}
$S=8-8=0$，常数区输出0。中心70、其余8个邻居50时，
\[
 g=8\cdot70-8\cdot50=560-400=160.
\]
中心越偏离周围，绝对响应越大。\textbf{结论：}八邻域拉普拉斯类纯高通／二阶细节提取核，无单一选定朝向，包含对角比较；它本身不保留背景。需要最终锐化时还应规定如何与原图相加及使用多大系数。

\subsubsection{A2③：沿行方向二阶差分，同时沿列方向求和}
该核可写成竖直系数列与水平系数行的外积：
\[
 \begin{bmatrix}-1\\2\\-1\end{bmatrix}
 \begin{bmatrix}1&1&1\end{bmatrix}
 =\begin{bmatrix}-1&-1&-1\\2&2&2\\-1&-1&-1\end{bmatrix}.
\]
即先把上、中、下三行各自的3个像素相加，分别记为$R_{\rm 上},R_{\rm 中},R_{\rm 下}$，再算
\[
 g=-R_{\rm 上}+2R_{\rm 中}-R_{\rm 下}.
\]
沿列方向的3点求和是一种\textbf{未除以3的平滑式汇集}；跨行的$[-1,2,-1]$则是负二阶差分。

\textbf{试验1：匀速渐变。}对三行分别全是10、20、30的窗口，
\[
 g=-3\cdot10+2(3\cdot20)-3\cdot30=-30+120-90=0.
\]
一阶变化存在，二阶变化却为0，已经说明它不是一般一阶Prewitt核。

\textbf{试验2：上下阶跃。}三行分别全是10、10、40时，
\[
 g=-30+60-120=-90;
\]
窗口下移，使三行变成10、40、40时，
\[
 g=-30+240-120=90.
\]
跨过同一个水平走向边界，产生一负一正的成对响应，符合二阶差分特征。

\textbf{试验3：左右阶跃。}若每一行都是$[10,10,40]$，则三行之和都为60：
\[
 g=-60+2\cdot60-60=0.
\]
\textbf{结论：}它是有方向性的零和二阶差分／高通类核，对\textbf{从上到下的灰度变化}敏感，常响应\textbf{水平走向边界}；对上述仅左右变化的理想图案为0。原讲义把它叫作“一阶差分／类Prewitt”不正确，应以系数$[-1,2,-1]$和数值试验为准。

\subsection{D1 原题与完整解答：中值滤波的五个问题}
\label{original:D1}
\textbf{题目【练习1原题】。}对中值滤波回答：①如何操作？②作用是什么？③是否可以用模板操作？④影响算法效率的因素有哪些？⑤试估计其频域特性。

\textbf{①如何操作。}选定窗口、锚点和边界规则；取出当前窗口数值；从小到大排序；奇数个样本取中间那个值，写到输出图对应位置；再移动窗口。9数窗口$[10,12,11,9,255,10,11,8,10]$排序后为$[8,9,10,10,10,11,11,12,255]$，第5个为10。符号写法为
\[
 g(r,c)=\operatorname{median}\{f(r+s,c+t):(s,t)\in W\},
\]
其中$W$就是窗口允许的偏移集合。若用偶数个样本，必须另行规定中间两个值如何处理，因此课堂通常用奇数尺寸。

\textbf{②作用。}非线性平滑，常用于抑制孤立脉冲和椒盐噪声。在合适条件下比同尺寸均值更少混合阶跃两侧的亮度；但会改变细线、尖角和小目标，也可能在高噪声比例时失效。

\textbf{③能否用模板。}能用\emph{窗口模板}选邻居；不能用\emph{固定线性系数模板}以卷积乘加实现所有输入的中值。可用$[0,0,1]$与$[0,1,0]$的加法反例证明其非线性，而不只背一句“因为要排序”。

\textbf{④效率因素。}设图像有$M\times N$个输出位置，窗口含$K$个数。朴素方法每个窗口完整排序，大致需要$O(K\log K)$次量级的比较，总量约$O(MNK\log K)$。这里$O(\cdot)$只表示随规模增长的数量级，不是精确运行秒数。窗口大小、图像尺寸、排序或选择算法、滑动时能否复用上一个窗口的信息、灰度是否为有限整数、内存访问与并行硬件都会影响速度。例如$3\times3$窗口需处理9个数，$7\times7$需处理49个数，不能认为只多做一两次运算。利用直方图更新等方法可明显改变复杂度。

\textbf{⑤频域特性。}中值是非线性的，不能给它一个对所有输入通用的固定传递函数$H$，再一律写“输出频谱=输入频谱乘$H$”。可以定性说它常抑制孤立的快速起伏，呈现某些平滑／低通效果；但其作用依赖信号形状、幅度、窗口和噪声。用正弦条纹测试时，非线性还可能产生新频率，所得只能是\textbf{该测试条件下的等效响应}，不是唯一线性频率响应。

\subsection{D2 原题与完整解答：均值和中值如何选择}
\label{original:D2}
\textbf{题目。}比较均值滤波与中值滤波在噪声抑制上的差异，分别适用哪种噪声？

\textbf{答题的出发点是计算规则。}均值把所有样本幅值都加入总和；一个很大的错误值会拉动平均。中值看排序位置，少数极端值通常落在两端，未必改变中间值。因此前面的含255窗口，均值为37.33，中值为10，差异不是算法“更聪明”，而是对异常值的敏感性不同。
\begin{center}
\begin{tabularx}{\linewidth}{@{}p{.16\linewidth}XX@{}}
\toprule
比较项目 & 均值滤波 & 中值滤波\\
\midrule
计算 & 求和后除以样本数；线性 & 排序取中间值；非线性\\
典型适用 & 局部平坦区域中的零均值、近似独立加性噪声，如高斯噪声 & 少数位置的椒盐／脉冲噪声\\
边缘代价 & 混合边缘两侧数值，形成过渡灰 & 常较少产生过渡灰，但会改动形状、删除小目标\\
极值影响 & 极值会拉偏均值，并随窗口移动扩散影响 & 少量极值通常不成为中值；比例过高仍会失败\\
\bottomrule
\end{tabularx}
\end{center}
\textbf{结论：}考试的基本匹配是“近似独立加性高斯噪声可考虑均值／高斯平滑，椒盐噪声优先考虑中值”。这不是不看图像结构的绝对排名。若要保留非常细的线，二者都可能有明显损失；若噪声强相关，独立噪声的方差下降公式也不能直接套用。

\subsection{D3 原题与完整解答：为何系数和是1或0}
\label{original:D3}
\textbf{题目。}为什么平滑（低通）模板的系数和等于1，而纯锐化（高通）模板的系数和等于0？

\textbf{先修正题目用语。}这里“系数和0”指\textbf{纯细节提取／零直流高通核}，不是所有最终锐化核。很多最终锐化核恰好和为1。

设常数输入为$C=100$。核系数为$w_1,\ldots,w_K$，输出为
\[
 g=\sum_jw_j\cdot100=100\sum_jw_j.
\]
希望平坦亮度仍为100，就令$\sum_jw_j=1$；希望完全没有变化的区域不产生细节，就令$\sum_jw_j=0$。举例：
\[
 \begin{array}{ll}
 \text{均值核：}&100\times(9/9)=100,\\
 \text{拉普拉斯细节核：}&100\times(4-4)=0,\\
 \text{原图+拉普拉斯：}&100\times(5-4)=100.
 \end{array}
\]
\textbf{补充条件：}图像边界采用保持常数的延拓规则，且未发生裁剪；“和为1”本身只保证常数响应，不保证整个核一定低通。对有限图像，要讨论全图平均值是否严格保持，还必须考虑边界实现，不能把局部常数测试扩大成任意边界下的全局平均定理。

\subsection{D4 原题与完整解答：拉普拉斯公式及一、二阶差异}
\label{original:D4}
\textbf{题目。}写出拉普拉斯锐化的公式与典型掩模，并说明一阶导数与二阶导数在锐化中的区别。

\textbf{第一部分：符号配套写全。}
\[
 \nabla^2f\ \approx\ Lf
 =f_{\rm 上}+f_{\rm 下}+f_{\rm 左}+f_{\rm 右}-4f_{\rm 中},
 \quad L=\begin{bmatrix}0&1&0\\1&-4&1\\0&1&0\end{bmatrix},
 \quad g=f-Lf.
\]
若使用中心为正4的$P=-L$，则写成
\[
 P=\begin{bmatrix}0&-1&0\\-1&4&-1\\0&-1&0\end{bmatrix},
 \qquad g=f+Pf.
\]
两套符号给出同一个最终核$[0,-1,0;-1,5,-1;0,-1,0]$。原讲义把“邻居和减4倍中心”的公式与“中心正4”的核并排对应，符号不一致；不能照抄后再任意选加或减。

\textbf{数值验算。}中心70、邻居50：$Lf=200-280=-80$，于是$g=70-(-80)=150$；或$Pf=80$，$g=70+80=150$。假如错误地对正中心核采用$70-80=-10$，局部亮峰被反向压低，说明符号配错。

\textbf{第二部分：一阶与二阶。}一阶问相邻像素变化多少；二阶问变化量怎样变化。前述0到10阶跃的一阶有一个$+10$主要响应，二阶在阶跃两侧出现$+10,-10$；匀速斜坡一阶非零、二阶为0。二阶对细线、孤立点及噪声敏感，常用于细节增强；梯度常用于检测轮廓与方向。实际边缘宽度还取决于核、尺度和阈值，不能把“粗／细”作为绝对定义。

\subsection{D5 原题与完整解答：非锐化掩模和高提升的关系}
\label{original:D5}
\textbf{题目。}说明非锐化掩模与高提升滤波的关系。

\textbf{先定义中间结果：}$m=f-f_{\rm blur}$是原图减模糊图得到的细节掩模。它在常数区为0，单独显示时不再保留原图背景；显示正负细节还需要选择灰度映射。

\textbf{再定义最终输出：}常见非锐化掩模锐化为$g=f+km$，$k=1$时$g=2f-f_{\rm blur}$。增加$k$可加强细节，但也会放大噪声和明暗边。

\textbf{与原题高提升写法对照：}
\[
 g_A=Af-f_{\rm blur}=(A-1)f+m.
\]
$A=1$时是纯掩模$m$；$A=2$时等于$k=1$的保常数锐化。一般$A>1$的常数增益为$A-1$，不能无条件说“保持原图亮度”。若要归一化，应写
\[
 g_{\rm norm}=\frac{Af-f_{\rm blur}}{A-1}
 =f+\frac{m}{A-1}.
\]
\textbf{数值检查：}原图80、模糊图80，在任何$k$下$f+km=80$；而$A=3$时$Af-f_{\rm blur}=240-80=160$。原图100、模糊图90时，掩模10；$k=1$输出110，$A=2$也输出110；$A=3$输出210，若除以2归一化则为105。先检查参数约定，再讨论锐化程度。

\subsection{把计算结果变成能交卷的解释}
遇到陌生核，可以按下面的次序写答案，而不必靠外形猜名称：
\begin{enumerate}[label=\arabic*.]
 \item \textbf{计算系数和。}说明常数输入$C$会变成什么；这一步区分是否保留背景。
 \item \textbf{检查正负和位置。}全非负且归一化的局部平均具有平滑意义；有正负项时，说明是在比较哪些位置。
 \item \textbf{选择一个小图案手算。}常数、孤立亮点、上下阶跃、左右阶跃或斜线都可以。一次只比较一个因素。
 \item \textbf{说清输出是哪一层。}是噪声较少的图、纯细节图，还是已经把原图加回的锐化图？
 \item \textbf{补上限制。}边界规则、取整裁剪、噪声放大、细线损失，以及方向性不是绝对各向同性等。
\end{enumerate}
本章最值得记住的不是核的外号，而是三个可检验的问题：\textbf{每个邻居怎样参与计算？一块常数会输出什么？一个明确的变化图案会输出什么？}能回答这三个问题，就能把大多数滤波判断题变成清楚的算术。
